\section*{Chapitre \ref{ch:fm:the-head-of-desks-job} --- Le travail du responsable de desk}

\begin{solution}{exo:fm:the-head-of-desks-job:1}
$\Phi(1\times(1-1.1))=\Phi(-0.1)=46.0~\%$.
\end{solution}

\begin{solution}{exo:fm:the-head-of-desks-job:2}
$\Phi(-(55-28)/42.2)=\Phi(-0.64)=26~\%$.
\end{solution}

\begin{solution}{exo:fm:the-head-of-desks-job:3}
$(55+20)/1.645=45.6$~; VaR à 99~\% quotidienne $2.326\times45.6/\sqrt{252}=6.7$.
\end{solution}

\begin{solution}{exo:fm:the-head-of-desks-job:4}
La mise à l'échelle exige $\mu-z_q\sigma>0$, c'est-à-dire $S>z_{0.95}=1.645$~; le 1.3 du desk n'est pas suffisant, donc aucune échelle n'atteint 95~\% pour un budget positif.
\end{solution}

\begin{solution}{exo:fm:the-head-of-desks-job:5}
Ratio de Sharpe de 0.5 pour $k=1.5$~: $\Phi(-0.25)=40.1~\%$. Un ratio de Sharpe de 2 exige $2(1-k)=-0.25$~: $k=1.125$.
\end{solution}

\begin{solution}{exo:fm:the-head-of-desks-job:6}
Budget cumulé à ce jour $60.5\times3/12=15.125$ contre $55\times3/12=13.75$~: 1.4~millions de dollars de retard. L'écart type sur trois mois est
$42.2\sqrt{0.25}=21.1$~: 0.065 écarts types, du bruit.
\end{solution}

\begin{solution}{exo:fm:the-head-of-desks-job:7}
Budget de 66~: $\lambda=66/(55-0.674\times42.2)=2.49$~; VaR à 99~\% quotidienne $2.326\times2.49\times42.2/\sqrt{252}=15.4$, contre 6.7 autorisés.
\end{solution}

\begin{solution}{exo:fm:the-head-of-desks-job:8}
Avec un budget supérieur de 10~\% à l'espérance, le desk manque son objectif avec une probabilité de 55~\% par an, et deux années de suite avec une probabilité de 30.5~\% si les années sont
indépendantes. Deux échecs disent peu de chose~; il faut comparer les revenus à l'espérance et le ratio de Sharpe à son historique.
\end{solution}

\begin{solution}{pb:fm:the-head-of-desks-job:1}
\begin{enumerate}
\item Voir la \cref{def:fm:the-head-of-desks-job:budget}.
\item 55, 42.2 et 1.30.
\item Environ 20.5~millions de dollars après frais et rémunération~; 26~\% de probabilité que les revenus soient inférieurs aux 28~millions de dollars de frais.
\item 60.5~millions de dollars.
\item Voir la \cref{prop:fm:the-head-of-desks-job:target}.
\item 44.8~\% sous forme fermée, 45.0~\% sur 20\,000 années simulées.
\item Ses résultats se regroupent autour de la moyenne, si bien qu'un objectif supérieur à la moyenne est plus éloigné en unités de sa volatilité.
\item 30.0~millions et 50.4~millions de dollars.
\item 2.28 fois chaque position.
\item Voir la \cref{def:fm:the-head-of-desks-job:appetite}~; la borne supérieure est la capacité de risque de la firme.
\item $(55+20)/1.645=45.6$.
\item 6.7 par jour~; seuils d'arrêt de 8.8, 7.0 et 4.2.
\item De 8~\% (45.6 contre 42.2).
\item L'\omterm{def:fm:the-head-of-desks-job:appetite}{appétit pour le risque} autorise 45.6 de volatilité~; atteindre le budget trois années sur quatre exige environ 96.
\item Voir la \cref{def:fm:the-head-of-desks-job:scope}~; sa valorisation, son enregistrement (booking) et son règlement-livraison sont nouveaux et non testés.
\item 1.4~millions de dollars, 0.065 écarts types.
\item Ajouter du risque pour rattraper~: l'écart est du bruit, et ajouter du risque après des pertes brise l'\omterm{def:fm:the-head-of-desks-job:appetite}{appétit pour le risque}.
\item Le budget et les limites de risque~; le ratio de Sharpe du desk ne relève pas d'une décision de décembre.
\item $\Phi(S(1-k))$~: 44.8~\% pour $S=1.3$, $k=1.1$~; une échelle de risque de 2.28 pour l'atteindre avec une probabilité de 75~\%.
\item Le budget est supérieur aux revenus espérés du desk et n'est donc atteint moins de la moitié du temps avec les limites actuelles~; soit on le fixe à l'espérance, soit
on relève les limites du facteur dont il a besoin, ce que l'\omterm{def:fm:the-head-of-desks-job:appetite}{appétit pour le risque} de la firme ne permet pas.
\end{enumerate}
\end{solution}

\begin{solution}{iq:fm:the-head-of-desks-job:1}
50~\%, et $\Phi(-1)=15.9~\%$ pour une perte.
\iqlookfor{quantiles normaux en unités de Sharpe.}
\end{solution}

\begin{solution}{iq:fm:the-head-of-desks-job:2}
Un énoncé écrit des pertes et des risques que la firme accepte, en résultat, en capital et en liquidité~; décliné en budgets de volatilité, limites de VaR et
seuils d'arrêt par desk et par métier.
\iqlookfor{la déclinaison descendante.}
\end{solution}

\begin{solution}{iq:fm:the-head-of-desks-job:3}
Avec un ratio de Sharpe de 1.3, la cote passe de 50~\% à $\Phi(-0.26)=39.7~\%$~; demander les limites dont le budget a besoin ou un budget à l'espérance.
\iqlookfor{$\Phi(S(1-k))$ et le lien entre budget et risque.}
\end{solution}

\begin{solution}{iq:fm:the-head-of-desks-job:4}
Quand l'écart est grand par rapport à la volatilité des revenus à ce jour, qui croît avec la racine carrée du temps~; au bout d'un trimestre, la plupart des écarts sont du bruit.
\iqlookfor{mise à l'échelle de la bande avec $\sqrt t$.}
\end{solution}

\begin{solution}{iq:fm:the-head-of-desks-job:5}
\omterm{def:fm:the-head-of-desks-job:scope}{Approbation des nouveaux produits}~: le risque, le product control, les opérations, le juridique, la fiscalité, la conformité et la technologie l'examinent et le valident, souvent avec des limites de taille initiales
limites~; ses contrôles sont les moins testés.
\iqlookfor{le processus d'approbation et pourquoi.}
\end{solution}

\begin{solution}{iq:fm:the-head-of-desks-job:6}
$P=\Phi(S(1-k))$ décroît en $S$ lorsque $k>1$ (un budget ambitieux) et croît lorsque $k<1$ (un budget prudent).
\iqlookfor{le signe de $1-k$.}
\end{solution}
